\documentclass[11pt,a4paper]{article}
\usepackage[margin=0.75in]{geometry}
\usepackage{graphicx}
\usepackage{hyperref}
\usepackage{tabularx}
\usepackage{booktabs}
\usepackage{float}
\usepackage{xcolor}

\title{\textbf{Task 4: GGUF Quantization \& Benchmarking}}
\author{TinyLlama MLX Fine-Tuning Project}
\date{\today}

\begin{document}

\maketitle

\section*{Task Description}
This task implements \textbf{post-training quantization} using the GGUF (GPT-Generated Unified Format) standard, converting the fine-tuned TinyLlama model to compressed formats suitable for CPU inference. We quantize to three precision levels and benchmark size/performance tradeoffs.

\textbf{Key objectives:}
\begin{itemize}
    \item Convert fine-tuned LoRA adapter to a fused HuggingFace model
    \item Quantize to GGUF formats: Q4\_K\_M (4-bit), Q5\_K\_M (5-bit), Q8\_0 (8-bit)
    \item Measure model size reduction and inference throughput using llama.cpp
\end{itemize}

\section*{Technical Implementation}

\subsection*{Procedure}
\begin{enumerate}
    \item \textbf{LoRA Fusion:} Merge LoRA adapters (rank-8) into base TinyLlama weights using MLX fuse API
    \item \textbf{Export to HuggingFace:} Save fused model in safetensors format with tokenizer and config
    \item \textbf{GGUF Conversion:} Use llama.cpp's \texttt{convert\_hf\_to\_gguf.py} to create FP16 GGUF baseline
    \item \textbf{Quantization:} Apply llama.cpp's \texttt{llama-quantize} with Q4\_K\_M, Q5\_K\_M, Q8\_0 methods
    \item \textbf{Benchmarking:} Measure prompt processing speed using llama.cpp's CLI tools
\end{enumerate}

\subsection*{Technical Details}
\begin{itemize}
    \item \textbf{Source Model:} TinyLlama + LoRA rank-8 adapter from Task 1 results
    \item \textbf{Tools:} llama.cpp (build b4378), MLX-LM fuse module
    \item \textbf{Quantization Methods:}
    \begin{itemize}
        \item Q4\_K\_M: 4-bit K-quants (mixed precision), balanced accuracy/size
        \item Q5\_K\_M: 5-bit K-quants, higher precision retention
        \item Q8\_0: 8-bit quantization, near-lossless quality
    \end{itemize}
    \item \textbf{Hardware:} Apple Silicon Mac (Metal GPU for MLX, CPU for llama.cpp)
\end{itemize}

\section*{Results}

\subsection*{Quantitative Metrics}
\begin{table}[H]
\centering
\scriptsize
\begin{tabular}{@{}lcccccc@{}}
\toprule
\textbf{Format} & \textbf{Size (MB)} & \textbf{Compression} & \textbf{Perplexity} & \textbf{PPL Loss} & \textbf{Tokens/sec} & \textbf{Memory (MB)} \\
\midrule
Original (FP16) & 2,099 & 1.00× & 4.2722 & 0.0000 & N/A & N/A \\
Q8\_0 & 1,116 & 1.88× & 4.2725 & 0.0003 & 157.1 & 1,224 \\
Q5\_K\_M & 746 & 2.81× & 4.2859 & 0.0137 & 166.4 & 862 \\
Q4\_K\_M & 637 & \textbf{3.30×} & 4.3399 & 0.0677 & \textbf{214.5} & \textbf{749} \\
\bottomrule
\end{tabular}
\caption{Model size, perplexity, throughput and memory across quantization formats}
\end{table}

\subsection*{Key Findings}
\begin{itemize}
    \item \textbf{3.30× compression:} Q4\_K\_M reduces model from 2.1 GB to 637 MB with 1.6\% perplexity loss
    \item \textbf{Best throughput:} Q4\_K\_M achieves 214.5 tokens/sec, 36\% faster than Q8\_0
    \item \textbf{Memory efficient:} Q4\_K\_M uses only 749 MB RAM vs 1,224 MB for Q8\_0
    \item \textbf{Near-lossless quality:} Q8\_0 has only 0.0003 perplexity degradation (0.007\% loss)
    \item \textbf{Batch scaling:} Q4\_K\_M maintains consistent throughput across batch sizes 1, 4, and 8
    \item \textbf{GGUF portability:} Quantized models can run on CPU-only systems via llama.cpp or Ollama
    \item \textbf{Production ready:} Packaged inference server enables immediate deployment
\end{itemize}

\subsection*{Implementation Evidence}
\begin{verbatim}
$ ls -lh results/task4/
-rw-r--r--  637M  model-Q4_K_M.gguf
-rw-r--r--  746M  model-Q5_K_M.gguf
-rw-r--r-- 1116M  model-Q8_0.gguf
-rw-r--r-- 2099M  model.gguf (FP16 baseline)

$ ls -l results/task4/inference_server/
server.py  README.md  requirements.txt
\end{verbatim}

\textbf{Benchmark Results (10 generations × 100 tokens):}
\begin{verbatim}
Q4_K_M: 214.5 tok/s, 0.47s latency, 749 MB memory
Q5_K_M: 166.4 tok/s, 0.60s latency, 862 MB memory
Q8_0:   157.1 tok/s, 0.64s latency, 1224 MB memory
\end{verbatim}

\subsection*{Batched Inference Performance}
\begin{table}[H]
\centering
\small
\begin{tabular}{@{}lccc@{}}
\toprule
\textbf{Quantization} & \textbf{Batch-1 (tok/s)} & \textbf{Batch-4 (tok/s)} & \textbf{Batch-8 (tok/s)} \\
\midrule
Q4\_K\_M & 197.1 & 198.9 & 193.1 \\
Q5\_K\_M & 144.1 & 139.2 & 116.5 \\
Q8\_0 & 95.9 & 80.9 & 80.6 \\
\bottomrule
\end{tabular}
\caption{Throughput across batch sizes (tokens per second)}
\end{table}

\subsection*{Deployment Scenarios}
\begin{itemize}
    \item \textbf{Mobile/Edge (RECOMMENDED):} Q4\_K\_M — 637 MB, 214 tok/s, 749 MB RAM, 1.6\% quality loss
    \item \textbf{Server/API:} Q5\_K\_M — 746 MB, 166 tok/s, 0.3\% quality loss, good balance
    \item \textbf{Archival/Research:} Q8\_0 — near-lossless (0.007\% perplexity loss)
    \item \textbf{Production Deployment:} Packaged server in \texttt{results/task4/inference\_server/}
\end{itemize}

\section*{Conclusion}
GGUF quantization enables \textbf{3.30× model compression} while maintaining instruction-following quality. The Q4\_K\_M format is the optimal choice for edge deployment: smallest size (637 MB), fastest inference (214 tok/s), lowest memory (749 MB), with acceptable 1.6\% perplexity degradation.

\textbf{Benchmarking confirms:}
\begin{itemize}
    \item Perplexity measured on 80-row Alpaca validation corpus (14,958 tokens)
    \item 10 warm generations per model (100 tokens each) for latency/memory
    \item Batched inference tests (batch sizes 1, 4, 8) show consistent Q4\_K\_M performance
    \item Q8\_0 provides near-lossless quality (0.0003 perplexity loss) for research use
\end{itemize}

This task demonstrates the complete MLX → HuggingFace → GGUF → llama.cpp pipeline, providing a production-ready deployment path with packaged inference server code.

\end{document}
